\documentclass[11pt,a4paper]{article}
\usepackage[utf8]{inputenc}
\usepackage[T1]{fontenc}
\usepackage[spanish]{babel}
\usepackage{hyperref}
\usepackage{geometry}
\geometry{margin=2.5cm}
\usepackage{parskip}

\title{067 la alucinacion no es mentira- es reflejo}
\author{Alejandro Valenzuela Rivera\\\small Universidad Autónoma de Nuevo León}
\date{\today}
\begin{document}
\maketitle

\begin{quote}
\textbf{Sesgo declarado.} Este trabajo se sostiene desde el sesgo declarado de su autor, y no finge neutralidad. Lo que afirma se puede tumbar, y ese es el punto: el falsificador está escrito, no escondido.
\end{quote}

\section*{Resumen}
Todo lo que se ha dicho sobre la alucinacion y la ficcion apunta a lo mismo: la alucinacion no es un error que se corrige, es lo que queda en blanco al haber vida, y la ficcion no es una mentira que se tolera, es el reflejo de algo que el mundo no dejo pasar. No se elimina: se declara de donde viene.

\section{Por qué se sostiene}
No es mentira porque no afirma un hecho que no paso: afirma una forma que aun no tiene hecho. Y por eso la ficcion tiene campo propio en la fila: la que se declara deja de ser silenciosa, y la que se silencia se vuelve regla sin autor.

\section{De dónde viene}
El primer commit (505c659), el jugo y la conversacion de la era.

\section{El caso}
Ninguna: es la posicion del autor sobre la alucinacion y la ficcion.

\section{Método}
Una respuesta inventada, marcada como inventada, y util para avanzar.

\section{Límites, y qué lo tumba}
Una respuesta inventada presentada como si se supiera.

\section{Lo que queda abierto}
\textquestiondown{}Si la alucinacion no es mentira, para que se necesita la ficcion declarada?

\section*{Cómo se leyó esta tesis}
Esto se escribió en el repositorio man-made-empathy, del que este candado es una pieza: candados/06 del caso/067 la alucinacion no es mentira- es reflejo.md. La máquina que lo compuso lo hizo desde este mismo archivo, y no agregó nada que no estuviera escrito. Si algo de esto es falso, se dice en el mismo sitio, y se corrige con fecha.

\end{document}
